\section{Ιεραρχικός Αξιολογητής 3 Επιπέδων ($\epsilon$-Bypass)}
\label{sec:arch_multi_tier}

Ο πυρήνας της επιτάχυνσης βασίζεται στον αγωγό αξιολόγησης τριών επιπέδων (Multi-Tier $\epsilon$-Bypass Evaluator). Για κάθε υποψήφιο άτομο $\bm{x}^*$ του πληθυσμού, υπολογίζεται η Ευκλείδεια απόσταση $d_{\min} = \min_{\bm{x}_i \in \mathcal{D}} \|\bm{x}^* - \bm{x}_i\|$ προς το σύνολο των ήδη γνωστών σχεδίων:

\begin{equation}
\label{eq:tier_logic}
\text{Evaluation Tier}(\bm{x}^*) =
\begin{cases}
\text{\textbf{Tier 1 (Hit):}} & \text{εάν } d_{\min} < \epsilon_{\text{exact}} \implies \text{Επιστροφή } y_{\text{NN}} \\
\text{\textbf{Tier 2 (Surrogate):}} & \text{εάν } \epsilon_{\text{exact}} \le d_{\min} \le R \implies \text{Πρόβλεψη } \hat{y}_{\text{MLP}}(\bm{x}^*) \\
\text{\textbf{Tier 3 (Simulation):}} & \text{εάν } d_{\min} > R \implies \text{Εκτέλεση VLM σε Worker}
\end{cases}
\end{equation}

Η λογική αυτή εφαρμόζεται σε κάθε βήμα αξιολόγησης και οι τρεις κλάδοι αντιστοιχούν σε αυξανόμενο υπολογιστικό κόστος. Στη συνέχεια αναλύεται η λειτουργία κάθε Tier λεπτομερώς.

\subsection*{Tier 1 — Ακριβής Αντιστοίχιση στο LEAD DHT}

Ο πρώτος κλάδος της εξίσωσης~(\ref{eq:tier_logic}) ενεργοποιείται όταν $d_{\min} < \epsilon_{\text{exact}}$, όπου $\epsilon_{\text{exact}}$ είναι μια μικρή σταθερά ανοχής (τυπικά $\epsilon_{\text{exact}} = 10^{-4}$ στον κανονικοποιημένο χώρο παραμέτρων). Σε αυτή την περίπτωση, το άτομο $\bm{x}^*$ θεωρείται πρακτικώς ταυτόσημο με ήδη αξιολογημένο σχέδιο $\bm{x}_i \in \mathcal{D}$, και επιστρέφεται άμεσα η αποθηκευμένη τιμή $y_{\text{NN}}$ χωρίς καμία πρόσθετη υπολογιστική επιβάρυνση. Η ανάκτηση εκτελείται μέσω της gRPC κλήσης \texttt{KNNSearch} στον LEAD DHT, η οποία αξιοποιεί τον μαθημένο δείκτη RMI (Recursive Model Index) για προσεγγιστική εύρεση \cite{wang2025lead}. Η αποδοτικότητα του Tier~1 εξαρτάται κρίσιμα από τον ρυθμό επαναλαμβανόμενων γονιδίων στον πληθυσμό: στα πειράματα αξιολόγησης, το Tier~1 αντιστοιχούσε κατά μέσο όρο σε περίπου $28{,}3\%$ των αξιολογήσεων, αντικατοπτρίζοντας τον υψηλό βαθμό επαναχρησιμοποίησης γονιδίων σε προχωρημένες γενεές.

\subsection*{Tier 2 — Πρόβλεψη μέσω Surrogate με Πύλη Αβεβαιότητας}

Ο δεύτερος κλάδος ενεργοποιείται όταν $\epsilon_{\text{exact}} \le d_{\min} \le R$. Στην περίπτωση αυτή, το άτομο $\bm{x}^*$ βρίσκεται εντός της περιοχής επιρροής του εκπαιδευμένου μοντέλου surrogate, αλλά δεν ταυτίζεται με κανένα γνωστό σχέδιο. Ο Orchestrator υποβάλλει κλήση \texttt{Predict} στον Surrogate Node, ο οποίος επιστρέφει τόσο την προβλεπόμενη τιμή $\hat{y}_{\text{MLP}}(\bm{x}^*)$ όσο και την τυπική απόκλιση αβεβαιότητας $\sigma(\bm{x}^*)$. Η αποδοχή της πρόβλεψης εξαρτάται από μια \emph{πύλη εμπιστοσύνης}:
\begin{equation}
\label{eq:surrogate_gate}
\text{Αποδοχή Tier 2} \iff \sigma(\bm{x}^*) < \theta,
\end{equation}
όπου $\theta$ είναι το όριο αβεβαιότητας που βαθμονομείται πειραματικά. Εάν $\sigma(\bm{x}^*) \ge \theta$, η αξιολόγηση ανακατευθύνεται στο Tier~3 παρά το γεγονός ότι $d_{\min} \le R$. Η χρήση της τυπικής απόκλισης ως κριτηρίου αποδοχής αντιστοιχεί στη φιλοσοφία των surrogate-assisted εξελικτικών αλγορίθμων όπου ο έλεγχος ποιότητας προβλέψεων είναι κρίσιμος για την αποφυγή σφαλμάτων λόγω εξωπολλαπλασιασμού \cite{jin2002framework}. Στα πειράματα, το Tier~2 αντιστοιχούσε σε $47{,}2\%$ των αξιολογήσεων, αποφεύγοντας πλήρεις προσομοιώσεις VLM.

\subsection*{Tier 3 — Πλήρης Προσομοίωση μέσω Envoy και Workers}

Ο τρίτος κλάδος ενεργοποιείται όταν $d_{\min} > R$ ή όταν η πύλη αβεβαιότητας αποτύχει. Ο Orchestrator αποστέλλει gRPC κλήση \texttt{Evaluate} στον Envoy (θύρα 50055), ο οποίος δρομολογεί το αίτημα σε έναν διαθέσιμο Worker της συστοιχίας αξιοποιώντας αλγόριθμο round-robin με λαμβανόμενη υγεία υπηρεσίας (health-aware load balancing) \cite{klein2017envoy}. Ο Worker εκτελεί την πλήρη τρισδιάστατη επίλυση VLM μέσω AeroSandbox και επιστρέφει τον λόγο $L/D$. Το αποτέλεσμα εγγράφεται αμέσως στο LEAD DHT (\texttt{Put}) και διοχετεύεται στη sliding window buffer του Surrogate Node.

\subsection*{Online Retraining και Λόγος Παράκαμψης}

Κάθε $N_{\text{retrain}}$ νέα αποτελέσματα Tier~3 που προστίθενται στο κυλιόμενο παράθυρο εκκινούν αυτόματα μια ασύγχρονη επανεκπαίδευση του surrogate μοντέλου στο παρασκήνιο, ώστε το μοντέλο να προσαρμόζεται στις πλέον πρόσφατα εξερευνημένες περιοχές του χώρου σχεδίασης. Στα πειράματα χρησιμοποιήθηκε $N_{\text{retrain}} = 50$.

Ο συνολικός \emph{λόγος παράκαμψης} $\beta$ ορίζεται ως το ποσοστό αξιολογήσεων που αποφεύγουν την πλήρη προσομοίωση:
\begin{equation}
\label{eq:bypass_ratio}
\beta = \frac{N_1 + N_2}{N_{\text{total}}},
\end{equation}
όπου $N_1$, $N_2$ είναι ο αριθμός αξιολογήσεων που εξυπηρετήθηκαν από Tier~1 και Tier~2 αντίστοιχα, και $N_{\text{total}}$ ο συνολικός αριθμός αξιολογήσεων. Στα πειράματα αξιολόγησης επιτεύχθηκε $\beta = 75{,}5\%$, επιβεβαιώνοντας την αποδοτικότητα του ιεραρχικού αγωγού.

Τα αποτελέσματα του Tier 3 διοχετεύονται αυτόματα στο LEAD DHT και τροφοδοτούν το κυλιόμενο παράθυρο (sliding window buffer) της υπηρεσίας Surrogate, ενεργοποιώντας online retraining στο παρασκήνιο.
